\sectionguard
\section{Luke with the Septuagint: The Visitation and the Ark}

The road from the toll-gate climbs into the hill country of Juda, and a
philological claim waits at the top. Philological claims can be tested.
It is said that St.~Luke composed the Visitation narrative on the pattern
of the ark's journey to Jerusalem: Mary goes up into the hill country as
David went up; Elizabeth's cry, her question, John's leap, and the three
months' stay answer, point for point, the second panel of the story told
in section~13. If that is true, it is true in Greek---in the vocabulary
of the Septuagint, whose idiom the infancy narrative everywhere speaks,
set against the Greek of Luke's own sentences. The words can be counted,
and they have been counted: every concordance claim below rests on a
lemma-level search of the whole Rahlfs corpus, run and adversarially
re-run, not on secondary lists.\endnote{The Septuagint concordance results in this section were
obtained by lemma-level search of the complete Rahlfs 1935 corpus in the
machine-readable edition of E.~Wong,
\url{https://github.com/eliranwong/LXX-Rahlfs-1935}, independently rebuilt
and re-run in an adversarial verification pass, with the decisive loci
confirmed against a second instrument (Blue Letter Bible LXX, e.g.
\url{https://www.blueletterbible.org/lxx/1ch/28/18/}); New Testament counts
against the SBL Greek New Testament,
\url{https://github.com/LogosBible/SBLGNT}. The machine text is a
derivative electronic edition; a spot-check against the printed
Rahlfs--Hanhart remains on the study's open queue (section~90).} One naming
note: the Greek Bible files the books of Samuel as 1--4 Kingdoms, and the
Douay-Rheims as 1--4 Kings, so 2~Kings~6 (=~2~Samuel~6) is the chapter in
question under either name.
Here are the two openings, side by side.

\begin{witnessbox}{2 Kings 6:2 (= 2 Samuel 6:2), Douay-Rheims}
And David arose and went, with all the people that were with him of the men
of Juda to fetch the ark of God, upon which the name of the Lord of hosts
is invoked, who sitteth over it upon the cherubims.
\end{witnessbox}

\begin{witnessbox}{Luke 1:39, Douay-Rheims}
And Mary rising up in those days, went into the hill country with haste
into a city of Juda.\endnote{Douay-Rheims (Challoner) quotations in this
section are from the drbo.org chapter pages: Luke~1 at
\url{https://drbo.org/chapter/49001.htm}; 2~Kings~6 at
\url{https://drbo.org/chapter/10006.htm}. The 2~Kings~6 verses quoted in
full in section~13 appear here only in the fragments the philology
requires.}
\end{witnessbox}

Set side by side, the two sentences already carry the shape of the claim:
someone rises and goes, into the country of Juda, toward a house that will
hold the holy for a season. But a shape is not a proof. What follows takes
the alleged echoes one at a time, each with its Greek skeleton and its
honest weight---the strong with their caveats, the weak named as weak, one
outright negative printed as a negative. Then the scholarly debate is
reported as a debate. What the Church has since done with this reading
belongs to the next three sections; here the question is only what stands
on the page.

% Drawing: the two processions overlaid — 2 Kings 6 beside Luke 1, with the
% honest weight of each correspondence, and the hill country inset.
\begin{sketchfig}
\begin{tikzpicture}[scale=0.92]
% headers
\node[slabel] at (2.3,9.6) {\textbf{The ark goes up}\\(2 Kings 6)};
\node[slabel] at (11.7,9.6) {\textbf{Mary goes up}\\(Luke 1)};
% left rail
\draw[sketch] (0.0,7.55) rectangle (4.6,8.95);
\node[stiny] at (2.3,8.25) {``And David arose and went\ldots\\of the men of Juda'' (6:2)};
\draw[sketch] (0.0,5.55) rectangle (4.6,6.95);
\node[stiny] at (2.3,6.25) {``How shall the ark of the Lord\\come to me?'' (6:9)};
\draw[sketch] (0.0,3.55) rectangle (4.6,4.95);
\node[stiny] at (2.3,4.25) {``king David leaping and\\dancing before the Lord''\\(6:16)};
\draw[sketch] (0.0,1.4) rectangle (4.6,2.95);
\node[stiny] at (2.3,2.17) {``the ark of the Lord abode\ldots\\three months: and the Lord\\blessed Obededom'' (6:11)};
% right rail
\draw[sketch] (9.4,7.55) rectangle (14.0,8.95);
\node[stiny] at (11.7,8.25) {``And Mary rising up\ldots\ went\\into the hill country\ldots\ into\\a city of Juda'' (1:39)};
\draw[sketch] (9.4,5.55) rectangle (14.0,6.95);
\node[stiny] at (11.7,6.25) {``whence is this to me, that\\the mother of my Lord should\\come to me?'' (1:43)};
\draw[sketch] (9.4,3.55) rectangle (14.0,4.95);
\node[stiny] at (11.7,4.25) {``the infant leaped in her\\womb'' (1:41, 44)};
\draw[sketch] (9.4,1.4) rectangle (14.0,2.95);
\node[stiny] at (11.7,2.17) {``And Mary abode with her\\about three months'' (1:56)};
% links with weights, labels held clear above each arrow
\draw[slink] (4.75,8.05) -- (9.25,8.05);
\node[stiny] at (7.0,8.62) {the same rising-and-going:\\a common formula --- weight\\only inside the cluster};
\draw[slink] (4.75,6.02) -- (9.25,6.02);
\node[stiny] at (7.0,6.5) {the awed question:\\a partial verbal match};
\draw[slink] (4.75,4.02) -- (9.25,4.02);
\node[stiny] at (7.0,4.56) {motif, not vocabulary:\\Luke's leap-word is\\Rebekah's (Gen 25:22)};
\draw[slink] (4.75,2.1) -- (9.25,2.1);
\node[stiny] at (7.0,2.6) {\emph{menas treis} ---\\word for word};
% hill country inset
\node[stiny] at (7.0,0.72) {the hill country of Juda --- Luke names no town};
% mini-map: the one patch of ground both stories cross
\draw[pline] (4.55,-2.05) rectangle (9.45,0.15);
\draw[line width=0.11pt, black!41] (6.411,-1.972) -- (6.332,-1.962);
\draw[line width=0.12pt, black!47] (6.697,-1.960) -- (6.599,-1.948);
\draw[line width=0.14pt, black!52] (7.068,-1.952) -- (7.148,-1.964);
\draw[line width=0.15pt, black!57] (7.254,-1.966) -- (7.347,-1.981);
\draw[line width=0.14pt, black!54] (8.235,-1.946) -- (8.122,-1.965);
\draw[line width=0.14pt, black!53] (8.318,-1.960) -- (8.206,-1.979);
\draw[line width=0.13pt, black!51] (8.430,-1.973) -- (8.323,-1.991);
\draw[line width=0.11pt, black!44] (8.620,-1.951) -- (8.533,-1.966);
\draw[line width=0.13pt, black!48] (8.909,-1.980) -- (8.986,-1.967);
\draw[line width=0.14pt, black!53] (8.997,-1.948) -- (9.088,-1.933);
\draw[line width=0.15pt, black!57] (9.111,-1.976) -- (9.214,-1.959);
\draw[line width=0.16pt, black!61] (9.264,-1.985) -- (9.377,-1.966);
\draw[line width=0.11pt, black!40] (6.395,-1.855) -- (6.319,-1.845);
\draw[line width=0.11pt, black!42] (6.467,-1.843) -- (6.384,-1.832);
\draw[line width=0.12pt, black!47] (6.620,-1.890) -- (6.524,-1.878);
\draw[line width=0.12pt, black!48] (6.689,-1.837) -- (6.592,-1.825);
\draw[line width=0.11pt, black!42] (6.850,-1.912) -- (6.767,-1.902);
\draw[line width=0.14pt, black!51] (7.064,-1.866) -- (7.141,-1.878);
\draw[line width=0.15pt, black!54] (7.116,-1.872) -- (7.202,-1.885);
\draw[line width=0.15pt, black!56] (7.289,-1.920) -- (7.380,-1.936);
\draw[line width=0.14pt, black!53] (7.396,-1.861) -- (7.476,-1.880);
\draw[line width=0.11pt, black!41] (7.714,-1.892) -- (7.637,-1.911);
\draw[line width=0.12pt, black!45] (7.844,-1.906) -- (7.753,-1.922);
\draw[line width=0.12pt, black!47] (7.891,-1.864) -- (7.795,-1.881);
\draw[line width=0.13pt, black!52] (8.072,-1.853) -- (7.963,-1.872);
\draw[line width=0.14pt, black!53] (8.167,-1.900) -- (8.054,-1.919);
\draw[line width=0.14pt, black!54] (8.212,-1.860) -- (8.099,-1.879);
\draw[line width=0.13pt, black!52] (8.392,-1.860) -- (8.284,-1.878);
\draw[line width=0.13pt, black!51] (8.446,-1.915) -- (8.340,-1.933);
\draw[line width=0.11pt, black!43] (8.620,-1.868) -- (8.535,-1.882);
\draw[line width=0.11pt, black!40] (8.668,-1.860) -- (8.591,-1.873);
\draw[line width=0.13pt, black!50] (8.924,-1.838) -- (9.007,-1.824);
\draw[line width=0.15pt, black!55] (9.024,-1.859) -- (9.120,-1.843);
\draw[line width=0.15pt, black!58] (9.105,-1.875) -- (9.209,-1.858);
\draw[line width=0.16pt, black!61] (9.239,-1.848) -- (9.351,-1.829);
\draw[line width=0.11pt, black!40] (6.409,-1.785) -- (6.332,-1.775);
\draw[line width=0.11pt, black!42] (6.467,-1.788) -- (6.384,-1.778);
\draw[line width=0.12pt, black!47] (6.619,-1.736) -- (6.524,-1.724);
\draw[line width=0.12pt, black!48] (6.687,-1.750) -- (6.589,-1.738);
\draw[line width=0.12pt, black!45] (6.804,-1.751) -- (6.712,-1.740);
\draw[line width=0.11pt, black!40] (6.898,-1.802) -- (6.822,-1.793);
\draw[line width=0.15pt, black!56] (7.186,-1.771) -- (7.277,-1.785);
\draw[line width=0.15pt, black!56] (7.253,-1.784) -- (7.345,-1.800);
\draw[line width=0.15pt, black!55] (7.330,-1.745) -- (7.418,-1.762);
\draw[line width=0.14pt, black!51] (7.439,-1.801) -- (7.513,-1.824);
\draw[line width=0.11pt, black!42] (7.724,-1.773) -- (7.644,-1.792);
\draw[line width=0.12pt, black!44] (7.786,-1.780) -- (7.699,-1.797);
\draw[line width=0.13pt, black!49] (7.930,-1.754) -- (7.830,-1.771);
\draw[line width=0.13pt, black!52] (8.058,-1.729) -- (7.949,-1.747);
\draw[line width=0.14pt, black!53] (8.140,-1.797) -- (8.028,-1.815);
\draw[line width=0.14pt, black!54] (8.234,-1.747) -- (8.121,-1.766);
\draw[line width=0.14pt, black!52] (8.362,-1.754) -- (8.253,-1.772);
\draw[line width=0.13pt, black!49] (8.465,-1.758) -- (8.364,-1.775);
\draw[line width=0.12pt, black!44] (8.580,-1.759) -- (8.492,-1.774);
\draw[line width=0.13pt, black!48] (8.892,-1.805) -- (8.971,-1.792);
\draw[line width=0.15pt, black!54] (8.996,-1.749) -- (9.091,-1.733);
\draw[line width=0.16pt, black!59] (9.129,-1.730) -- (9.237,-1.712);
\draw[line width=0.16pt, black!61] (9.281,-1.801) -- (9.394,-1.782);
\draw[line width=0.12pt, black!44] (5.309,-1.653) -- (5.271,-1.723);
\draw[line width=0.12pt, black!46] (5.356,-1.689) -- (5.329,-1.762);
\draw[line width=0.14pt, black!51] (5.458,-1.663) -- (5.449,-1.746);
\draw[line width=0.15pt, black!55] (5.591,-1.628) -- (5.614,-1.712);
\draw[line width=0.15pt, black!54] (5.708,-1.623) -- (5.752,-1.692);
\draw[line width=0.11pt, black!43] (6.500,-1.673) -- (6.416,-1.663);
\draw[line width=0.12pt, black!46] (6.584,-1.698) -- (6.492,-1.687);
\draw[line width=0.12pt, black!47] (6.747,-1.637) -- (6.650,-1.625);
\draw[line width=0.11pt, black!44] (6.857,-1.652) -- (6.771,-1.642);
\draw[line width=0.14pt, black!53] (7.135,-1.623) -- (7.218,-1.636);
\draw[line width=0.15pt, black!56] (7.230,-1.696) -- (7.321,-1.711);
\draw[line width=0.15pt, black!54] (7.381,-1.658) -- (7.464,-1.677);
\draw[line width=0.11pt, black!42] (7.724,-1.685) -- (7.644,-1.705);
\draw[line width=0.12pt, black!47] (7.851,-1.622) -- (7.756,-1.639);
\draw[line width=0.13pt, black!48] (7.891,-1.649) -- (7.793,-1.666);
\draw[line width=0.13pt, black!52] (8.041,-1.666) -- (7.932,-1.685);
\draw[line width=0.14pt, black!53] (8.102,-1.699) -- (7.991,-1.718);
\draw[line width=0.14pt, black!54] (8.220,-1.661) -- (8.107,-1.680);
\draw[line width=0.14pt, black!52] (8.326,-1.621) -- (8.216,-1.639);
\draw[line width=0.12pt, black!47] (8.495,-1.633) -- (8.400,-1.649);
\draw[line width=0.12pt, black!45] (8.547,-1.637) -- (8.458,-1.652);
\draw[line width=0.13pt, black!49] (8.869,-1.647) -- (8.949,-1.634);
\draw[line width=0.15pt, black!54] (8.981,-1.680) -- (9.075,-1.664);
\draw[line width=0.16pt, black!60] (9.133,-1.658) -- (9.242,-1.639);
\draw[line width=0.16pt, black!61] (9.262,-1.702) -- (9.375,-1.683);
\draw[line width=0.11pt, black!42] (5.190,-1.513) -- (5.121,-1.560);
\draw[line width=0.12pt, black!44] (5.257,-1.520) -- (5.191,-1.580);
\draw[line width=0.13pt, black!49] (5.352,-1.561) -- (5.307,-1.641);
\draw[line width=0.15pt, black!57] (5.534,-1.541) -- (5.547,-1.636);
\draw[line width=0.16pt, black!58] (5.576,-1.538) -- (5.605,-1.630);
\draw[line width=0.16pt, black!57] (5.733,-1.535) -- (5.794,-1.600);
\draw[line width=0.11pt, black!42] (6.480,-1.594) -- (6.399,-1.584);
\draw[line width=0.12pt, black!45] (6.577,-1.510) -- (6.487,-1.501);
\draw[line width=0.12pt, black!48] (6.734,-1.521) -- (6.636,-1.512);
\draw[line width=0.12pt, black!46] (6.813,-1.558) -- (6.719,-1.548);
\draw[line width=0.15pt, black!54] (7.176,-1.577) -- (7.264,-1.589);
\draw[line width=0.15pt, black!56] (7.272,-1.558) -- (7.363,-1.571);
\draw[line width=0.15pt, black!55] (7.345,-1.544) -- (7.433,-1.558);
\draw[line width=0.11pt, black!41] (7.706,-1.544) -- (7.627,-1.563);
\draw[line width=0.12pt, black!47] (7.847,-1.507) -- (7.751,-1.520);
\draw[line width=0.13pt, black!50] (7.945,-1.541) -- (7.841,-1.559);
\draw[line width=0.14pt, black!53] (8.068,-1.535) -- (7.957,-1.554);
\draw[line width=0.14pt, black!53] (8.150,-1.592) -- (8.037,-1.611);
\draw[line width=0.14pt, black!53] (8.219,-1.533) -- (8.106,-1.552);
\draw[line width=0.13pt, black!50] (8.402,-1.509) -- (8.299,-1.527);
\draw[line width=0.12pt, black!47] (8.470,-1.510) -- (8.374,-1.526);
\draw[line width=0.11pt, black!43] (8.557,-1.509) -- (8.472,-1.524);
\draw[line width=0.14pt, black!52] (8.913,-1.519) -- (9.002,-1.504);
\draw[line width=0.15pt, black!57] (9.020,-1.561) -- (9.121,-1.544);
\draw[line width=0.16pt, black!60] (9.160,-1.587) -- (9.271,-1.569);
\draw[line width=0.16pt, black!61] (9.248,-1.538) -- (9.361,-1.519);
\draw[line width=0.11pt, black!42] (5.159,-1.460) -- (5.084,-1.495);
\draw[line width=0.12pt, black!46] (5.268,-1.429) -- (5.185,-1.475);
\draw[line width=0.12pt, black!46] (5.405,-1.400) -- (5.342,-1.466);
\draw[line width=0.15pt, black!56] (5.537,-1.413) -- (5.564,-1.498);
\draw[line width=0.16pt, black!58] (5.566,-1.431) -- (5.608,-1.514);
\draw[line width=0.16pt, black!60] (5.694,-1.428) -- (5.763,-1.497);
\draw[line width=0.15pt, black!56] (5.858,-1.427) -- (5.920,-1.489);
\draw[line width=0.15pt, black!53] (5.946,-1.440) -- (6.001,-1.495);
\draw[line width=0.11pt, black!43] (6.527,-1.431) -- (6.442,-1.423);
\draw[line width=0.12pt, black!45] (6.595,-1.453) -- (6.504,-1.444);
\draw[line width=0.12pt, black!48] (6.711,-1.396) -- (6.614,-1.386);
\draw[line width=0.12pt, black!46] (6.830,-1.417) -- (6.736,-1.408);
\draw[line width=0.11pt, black!40] (6.928,-1.481) -- (6.849,-1.474);
\draw[line width=0.14pt, black!51] (7.140,-1.405) -- (7.221,-1.415);
\draw[line width=0.15pt, black!55] (7.227,-1.431) -- (7.317,-1.443);
\draw[line width=0.15pt, black!53] (7.382,-1.460) -- (7.466,-1.474);
\draw[line width=0.11pt, black!40] (7.676,-1.411) -- (7.600,-1.425);
\draw[line width=0.12pt, black!45] (7.785,-1.418) -- (7.694,-1.428);
\draw[line width=0.13pt, black!51] (7.963,-1.398) -- (7.856,-1.407);
\draw[line width=0.13pt, black!52] (8.003,-1.420) -- (7.893,-1.430);
\draw[line width=0.14pt, black!53] (8.120,-1.410) -- (8.006,-1.420);
\draw[line width=0.14pt, black!53] (8.232,-1.475) -- (8.119,-1.491);
\draw[line width=0.13pt, black!49] (8.403,-1.436) -- (8.301,-1.449);
\draw[line width=0.12pt, black!47] (8.477,-1.482) -- (8.382,-1.498);
\draw[line width=0.14pt, black!52] (8.902,-1.483) -- (8.990,-1.468);
\draw[line width=0.16pt, black!58] (9.054,-1.432) -- (9.160,-1.414);
\draw[line width=0.16pt, black!61] (9.160,-1.434) -- (9.272,-1.415);
\draw[line width=0.16pt, black!61] (9.271,-1.465) -- (9.384,-1.446);
\draw[line width=0.12pt, black!45] (5.185,-1.343) -- (5.096,-1.356);
\draw[line width=0.12pt, black!47] (5.313,-1.298) -- (5.216,-1.297);
\draw[line width=0.12pt, black!46] (5.374,-1.373) -- (5.293,-1.419);
\draw[line width=0.15pt, black!55] (5.571,-1.343) -- (5.628,-1.402);
\draw[line width=0.16pt, black!59] (5.691,-1.310) -- (5.757,-1.376);
\draw[line width=0.16pt, black!59] (5.809,-1.366) -- (5.877,-1.433);
\draw[line width=0.16pt, black!58] (5.915,-1.320) -- (5.980,-1.384);
\draw[line width=0.15pt, black!55] (6.021,-1.316) -- (6.078,-1.375);
\draw[line width=0.11pt, black!40] (6.457,-1.363) -- (6.380,-1.357);
\draw[line width=0.12pt, black!45] (6.601,-1.327) -- (6.510,-1.318);
\draw[line width=0.12pt, black!47] (6.671,-1.357) -- (6.575,-1.347);
\draw[line width=0.12pt, black!46] (6.840,-1.346) -- (6.747,-1.337);
\draw[line width=0.11pt, black!42] (6.917,-1.340) -- (6.834,-1.333);
\draw[line width=0.14pt, black!50] (7.130,-1.311) -- (7.206,-1.322);
\draw[line width=0.15pt, black!54] (7.220,-1.355) -- (7.308,-1.367);
\draw[line width=0.15pt, black!55] (7.343,-1.335) -- (7.431,-1.349);
\draw[line width=0.11pt, black!42] (7.719,-1.350) -- (7.637,-1.362);
\draw[line width=0.12pt, black!44] (7.775,-1.309) -- (7.685,-1.320);
\draw[line width=0.13pt, black!49] (7.903,-1.288) -- (7.800,-1.298);
\draw[line width=0.14pt, black!52] (8.022,-1.311) -- (7.912,-1.320);
\draw[line width=0.14pt, black!54] (8.145,-1.288) -- (8.031,-1.297);
\draw[line width=0.14pt, black!52] (8.292,-1.311) -- (8.182,-1.320);
\draw[line width=0.13pt, black!52] (8.322,-1.361) -- (8.213,-1.370);
\draw[line width=0.12pt, black!48] (8.435,-1.312) -- (8.336,-1.320);
\draw[line width=0.13pt, black!48] (8.797,-1.338) -- (8.874,-1.331);
\draw[line width=0.15pt, black!54] (8.908,-1.320) -- (9.002,-1.313);
\draw[line width=0.16pt, black!59] (9.016,-1.316) -- (9.122,-1.307);
\draw[line width=0.16pt, black!61] (9.088,-1.313) -- (9.199,-1.304);
\draw[line width=0.16pt, black!61] (9.253,-1.348) -- (9.366,-1.331);
\draw[line width=0.11pt, black!42] (5.135,-1.242) -- (5.053,-1.229);
\draw[line width=0.12pt, black!46] (5.238,-1.260) -- (5.144,-1.246);
\draw[line width=0.12pt, black!46] (5.378,-1.214) -- (5.301,-1.160);
\draw[line width=0.11pt, black!41] (5.509,-1.180) -- (5.452,-1.122);
\draw[line width=0.16pt, black!60] (5.838,-1.250) -- (5.906,-1.319);
\draw[line width=0.16pt, black!60] (5.969,-1.189) -- (6.037,-1.257);
\draw[line width=0.16pt, black!58] (6.017,-1.232) -- (6.080,-1.296);
\draw[line width=0.11pt, black!41] (6.520,-1.230) -- (6.439,-1.223);
\draw[line width=0.12pt, black!45] (6.614,-1.262) -- (6.523,-1.253);
\draw[line width=0.12pt, black!47] (6.677,-1.208) -- (6.582,-1.199);
\draw[line width=0.12pt, black!47] (6.817,-1.226) -- (6.721,-1.217);
\draw[line width=0.12pt, black!44] (6.888,-1.234) -- (6.798,-1.226);
\draw[line width=0.14pt, black!52] (7.168,-1.229) -- (7.249,-1.240);
\draw[line width=0.15pt, black!54] (7.234,-1.260) -- (7.322,-1.272);
\draw[line width=0.15pt, black!55] (7.327,-1.226) -- (7.416,-1.240);
\draw[line width=0.11pt, black!43] (7.737,-1.251) -- (7.652,-1.263);
\draw[line width=0.12pt, black!44] (7.766,-1.181) -- (7.678,-1.192);
\draw[line width=0.13pt, black!49] (7.906,-1.193) -- (7.803,-1.203);
\draw[line width=0.14pt, black!53] (8.047,-1.225) -- (7.935,-1.235);
\draw[line width=0.14pt, black!53] (8.118,-1.242) -- (8.004,-1.251);
\draw[line width=0.14pt, black!53] (8.214,-1.261) -- (8.100,-1.271);
\draw[line width=0.13pt, black!51] (8.347,-1.185) -- (8.241,-1.194);
\draw[line width=0.12pt, black!44] (8.503,-1.212) -- (8.414,-1.220);
\draw[line width=0.11pt, black!40] (8.572,-1.241) -- (8.494,-1.248);
\draw[line width=0.15pt, black!54] (8.894,-1.186) -- (8.987,-1.178);
\draw[line width=0.16pt, black!59] (9.017,-1.221) -- (9.123,-1.212);
\draw[line width=0.16pt, black!61] (9.104,-1.220) -- (9.216,-1.211);
\draw[line width=0.16pt, black!62] (9.211,-1.251) -- (9.326,-1.242);
\draw[line width=0.11pt, black!42] (5.183,-1.114) -- (5.110,-1.072);
\draw[line width=0.12pt, black!46] (5.317,-1.099) -- (5.250,-1.032);
\draw[line width=0.12pt, black!47] (5.370,-1.116) -- (5.302,-1.047);
\draw[line width=0.12pt, black!47] (5.457,-1.105) -- (5.389,-1.037);
\draw[line width=0.16pt, black!58] (5.837,-1.152) -- (5.902,-1.217);
\draw[line width=0.16pt, black!59] (5.929,-1.071) -- (5.995,-1.137);
\draw[line width=0.16pt, black!59] (6.032,-1.145) -- (6.098,-1.213);
\draw[line width=0.16pt, black!58] (6.130,-1.085) -- (6.201,-1.143);
\draw[line width=0.11pt, black!41] (6.531,-1.145) -- (6.450,-1.138);
\draw[line width=0.12pt, black!45] (6.638,-1.132) -- (6.546,-1.123);
\draw[line width=0.12pt, black!47] (6.680,-1.140) -- (6.584,-1.130);
\draw[line width=0.12pt, black!47] (6.824,-1.087) -- (6.727,-1.078);
\draw[line width=0.11pt, black!43] (6.932,-1.131) -- (6.847,-1.124);
\draw[line width=0.14pt, black!51] (7.160,-1.109) -- (7.238,-1.121);
\draw[line width=0.15pt, black!53] (7.219,-1.127) -- (7.305,-1.139);
\draw[line width=0.15pt, black!54] (7.381,-1.094) -- (7.466,-1.110);
\draw[line width=0.14pt, black!52] (7.441,-1.113) -- (7.518,-1.131);
\draw[line width=0.12pt, black!47] (7.834,-1.138) -- (7.737,-1.148);
\draw[line width=0.13pt, black!50] (7.922,-1.093) -- (7.817,-1.103);
\draw[line width=0.14pt, black!53] (8.068,-1.085) -- (7.955,-1.095);
\draw[line width=0.14pt, black!54] (8.128,-1.071) -- (8.014,-1.081);
\draw[line width=0.14pt, black!52] (8.269,-1.108) -- (8.158,-1.118);
\draw[line width=0.13pt, black!51] (8.335,-1.091) -- (8.228,-1.100);
\draw[line width=0.12pt, black!45] (8.473,-1.100) -- (8.381,-1.108);
\draw[line width=0.14pt, black!52] (8.843,-1.083) -- (8.931,-1.076);
\draw[line width=0.15pt, black!56] (8.938,-1.117) -- (9.037,-1.109);
\draw[line width=0.15pt, black!58] (8.984,-1.138) -- (9.088,-1.129);
\draw[line width=0.16pt, black!61] (9.122,-1.068) -- (9.235,-1.059);
\draw[line width=0.16pt, black!62] (9.200,-1.111) -- (9.315,-1.102);
\draw[line width=0.11pt, black!40] (5.205,-1.007) -- (5.149,-0.953);
\draw[line width=0.11pt, black!43] (5.314,-0.980) -- (5.254,-0.920);
\draw[line width=0.12pt, black!46] (5.420,-0.995) -- (5.353,-0.928);
\draw[line width=0.12pt, black!46] (5.459,-0.961) -- (5.392,-0.894);
\draw[line width=0.12pt, black!47] (5.587,-0.976) -- (5.520,-0.909);
\draw[line width=0.11pt, black!41] (5.682,-1.018) -- (5.625,-0.962);
\draw[line width=0.15pt, black!54] (5.872,-1.038) -- (5.929,-1.095);
\draw[line width=0.15pt, black!53] (5.936,-0.962) -- (5.992,-1.018);
\draw[line width=0.16pt, black!59] (6.067,-0.990) -- (6.152,-1.036);
\draw[line width=0.15pt, black!57] (6.178,-0.996) -- (6.264,-1.027);
\draw[line width=0.14pt, black!53] (6.245,-0.965) -- (6.326,-0.981);
\draw[line width=0.11pt, black!43] (6.586,-0.983) -- (6.501,-0.975);
\draw[line width=0.12pt, black!47] (6.735,-1.039) -- (6.637,-1.029);
\draw[line width=0.12pt, black!47] (6.830,-0.958) -- (6.733,-0.949);
\draw[line width=0.11pt, black!43] (6.934,-0.985) -- (6.847,-0.977);
\draw[line width=0.14pt, black!51] (7.173,-0.988) -- (7.250,-1.000);
\draw[line width=0.15pt, black!54] (7.262,-0.958) -- (7.348,-0.971);
\draw[line width=0.15pt, black!55] (7.327,-1.002) -- (7.415,-1.017);
\draw[line width=0.11pt, black!40] (7.689,-0.983) -- (7.612,-0.999);
\draw[line width=0.12pt, black!46] (7.794,-0.961) -- (7.702,-0.971);
\draw[line width=0.13pt, black!50] (7.917,-0.963) -- (7.812,-0.973);
\draw[line width=0.14pt, black!53] (8.067,-0.957) -- (7.954,-0.967);
\draw[line width=0.14pt, black!53] (8.170,-0.978) -- (8.056,-0.988);
\draw[line width=0.14pt, black!52] (8.267,-1.010) -- (8.156,-1.019);
\draw[line width=0.13pt, black!51] (8.324,-1.020) -- (8.217,-1.029);
\draw[line width=0.12pt, black!45] (8.475,-1.015) -- (8.385,-1.022);
\draw[line width=0.14pt, black!52] (8.842,-0.993) -- (8.931,-0.986);
\draw[line width=0.15pt, black!57] (8.951,-0.959) -- (9.053,-0.950);
\draw[line width=0.16pt, black!60] (9.057,-0.994) -- (9.168,-0.984);
\draw[line width=0.16pt, black!62] (9.155,-0.977) -- (9.269,-0.967);
\draw[line width=0.16pt, black!62] (9.203,-1.029) -- (9.317,-1.019);
\draw[line width=0.12pt, black!44] (5.413,-0.930) -- (5.350,-0.866);
\draw[line width=0.12pt, black!46] (5.540,-0.853) -- (5.474,-0.787);
\draw[line width=0.12pt, black!47] (5.631,-0.925) -- (5.563,-0.857);
\draw[line width=0.12pt, black!47] (5.680,-0.858) -- (5.611,-0.790);
\draw[line width=0.15pt, black!57] (6.075,-0.923) -- (6.170,-0.935);
\draw[line width=0.14pt, black!54] (6.199,-0.896) -- (6.289,-0.895);
\draw[line width=0.14pt, black!50] (6.260,-0.886) -- (6.342,-0.882);
\draw[line width=0.12pt, black!44] (6.636,-0.909) -- (6.546,-0.900);
\draw[line width=0.12pt, black!47] (6.735,-0.880) -- (6.638,-0.871);
\draw[line width=0.12pt, black!47] (6.861,-0.900) -- (6.765,-0.891);
\draw[line width=0.12pt, black!45] (6.910,-0.913) -- (6.819,-0.904);
\draw[line width=0.11pt, black!40] (7.003,-0.873) -- (6.926,-0.867);
\draw[line width=0.14pt, black!53] (7.256,-0.848) -- (7.341,-0.861);
\draw[line width=0.15pt, black!54] (7.355,-0.883) -- (7.441,-0.899);
\draw[line width=0.14pt, black!51] (7.459,-0.925) -- (7.535,-0.945);
\draw[line width=0.11pt, black!42] (7.719,-0.926) -- (7.637,-0.939);
\draw[line width=0.13pt, black!48] (7.851,-0.887) -- (7.752,-0.897);
\draw[line width=0.13pt, black!50] (7.916,-0.870) -- (7.811,-0.879);
\draw[line width=0.13pt, black!52] (7.990,-0.891) -- (7.880,-0.900);
\draw[line width=0.14pt, black!53] (8.123,-0.857) -- (8.008,-0.867);
\draw[line width=0.14pt, black!52] (8.262,-0.930) -- (8.152,-0.939);
\draw[line width=0.13pt, black!51] (8.319,-0.906) -- (8.212,-0.915);
\draw[line width=0.11pt, black!43] (8.489,-0.846) -- (8.403,-0.853);
\draw[line width=0.11pt, black!40] (8.551,-0.919) -- (8.473,-0.926);
\draw[line width=0.14pt, black!51] (8.817,-0.912) -- (8.902,-0.905);
\draw[line width=0.15pt, black!55] (8.884,-0.932) -- (8.979,-0.924);
\draw[line width=0.16pt, black!59] (8.978,-0.897) -- (9.083,-0.889);
\draw[line width=0.16pt, black!61] (9.100,-0.875) -- (9.213,-0.866);
\draw[line width=0.16pt, black!61] (9.284,-0.913) -- (9.397,-0.904);
\draw[line width=0.11pt, black!40] (5.413,-0.801) -- (5.358,-0.746);
\draw[line width=0.11pt, black!43] (5.511,-0.784) -- (5.451,-0.723);
\draw[line width=0.12pt, black!46] (5.583,-0.810) -- (5.517,-0.744);
\draw[line width=0.12pt, black!47] (5.701,-0.804) -- (5.632,-0.735);
\draw[line width=0.11pt, black!43] (5.851,-0.819) -- (5.808,-0.746);
\draw[line width=0.11pt, black!42] (5.942,-0.817) -- (5.979,-0.742);
\draw[line width=0.12pt, black!47] (6.050,-0.747) -- (6.117,-0.677);
\draw[line width=0.13pt, black!49] (6.118,-0.754) -- (6.196,-0.700);
\draw[line width=0.12pt, black!46] (6.275,-0.805) -- (6.349,-0.781);
\draw[line width=0.12pt, black!44] (6.637,-0.764) -- (6.549,-0.755);
\draw[line width=0.12pt, black!46] (6.691,-0.798) -- (6.598,-0.789);
\draw[line width=0.12pt, black!48] (6.812,-0.797) -- (6.714,-0.787);
\draw[line width=0.11pt, black!43] (6.957,-0.816) -- (6.871,-0.808);
\draw[line width=0.15pt, black!54] (7.299,-0.754) -- (7.385,-0.768);
\draw[line width=0.15pt, black!54] (7.337,-0.822) -- (7.424,-0.836);
\draw[line width=0.14pt, black!52] (7.451,-0.764) -- (7.528,-0.784);
\draw[line width=0.12pt, black!48] (7.844,-0.804) -- (7.746,-0.815);
\draw[line width=0.13pt, black!51] (7.923,-0.740) -- (7.817,-0.749);
\draw[line width=0.14pt, black!53] (8.057,-0.814) -- (7.943,-0.823);
\draw[line width=0.14pt, black!53] (8.126,-0.822) -- (8.012,-0.831);
\draw[line width=0.14pt, black!52] (8.255,-0.770) -- (8.145,-0.779);
\draw[line width=0.13pt, black!49] (8.346,-0.759) -- (8.243,-0.768);
\draw[line width=0.12pt, black!46] (8.432,-0.757) -- (8.339,-0.765);
\draw[line width=0.14pt, black!51] (8.798,-0.800) -- (8.882,-0.793);
\draw[line width=0.15pt, black!58] (8.939,-0.768) -- (9.041,-0.759);
\draw[line width=0.16pt, black!60] (9.009,-0.762) -- (9.118,-0.753);
\draw[line width=0.16pt, black!62] (9.153,-0.804) -- (9.267,-0.794);
\draw[line width=0.16pt, black!62] (9.200,-0.810) -- (9.314,-0.800);
\draw[line width=0.12pt, black!45] (5.643,-0.705) -- (5.579,-0.642);
\draw[line width=0.12pt, black!45] (5.732,-0.654) -- (5.681,-0.578);
\draw[line width=0.12pt, black!47] (5.872,-0.700) -- (5.859,-0.604);
\draw[line width=0.12pt, black!47] (5.969,-0.689) -- (5.998,-0.597);
\draw[line width=0.12pt, black!44] (6.088,-0.636) -- (6.138,-0.563);
\draw[line width=0.11pt, black!43] (6.158,-0.661) -- (6.219,-0.601);
\draw[line width=0.11pt, black!41] (6.578,-0.639) -- (6.496,-0.630);
\draw[line width=0.12pt, black!46] (6.711,-0.684) -- (6.616,-0.675);
\draw[line width=0.12pt, black!48] (6.818,-0.703) -- (6.720,-0.694);
\draw[line width=0.11pt, black!44] (6.968,-0.659) -- (6.881,-0.652);
\draw[line width=0.14pt, black!53] (7.274,-0.676) -- (7.357,-0.690);
\draw[line width=0.15pt, black!54] (7.338,-0.663) -- (7.424,-0.679);
\draw[line width=0.11pt, black!43] (7.740,-0.640) -- (7.655,-0.652);
\draw[line width=0.12pt, black!46] (7.792,-0.665) -- (7.700,-0.676);
\draw[line width=0.13pt, black!50] (7.891,-0.707) -- (7.788,-0.717);
\draw[line width=0.14pt, black!53] (8.021,-0.631) -- (7.909,-0.641);
\draw[line width=0.14pt, black!53] (8.130,-0.639) -- (8.016,-0.649);
\draw[line width=0.13pt, black!52] (8.251,-0.702) -- (8.141,-0.711);
\draw[line width=0.13pt, black!50] (8.322,-0.692) -- (8.217,-0.701);
\draw[line width=0.11pt, black!42] (8.497,-0.706) -- (8.414,-0.713);
\draw[line width=0.14pt, black!53] (8.835,-0.680) -- (8.926,-0.672);
\draw[line width=0.15pt, black!57] (8.922,-0.692) -- (9.023,-0.684);
\draw[line width=0.16pt, black!61] (9.058,-0.713) -- (9.170,-0.704);
\draw[line width=0.16pt, black!62] (9.145,-0.630) -- (9.260,-0.621);
\draw[line width=0.16pt, black!62] (9.236,-0.671) -- (9.349,-0.661);
\draw[line width=0.11pt, black!41] (5.653,-0.603) -- (5.601,-0.540);
\draw[line width=0.11pt, black!41] (5.751,-0.533) -- (5.721,-0.458);
\draw[line width=0.12pt, black!44] (5.855,-0.569) -- (5.844,-0.481);
\draw[line width=0.11pt, black!43] (5.939,-0.548) -- (5.949,-0.462);
\draw[line width=0.11pt, black!43] (6.069,-0.581) -- (6.107,-0.506);
\draw[line width=0.11pt, black!41] (6.571,-0.551) -- (6.491,-0.541);
\draw[line width=0.12pt, black!47] (6.746,-0.521) -- (6.652,-0.501);
\draw[line width=0.12pt, black!48] (6.826,-0.604) -- (6.727,-0.594);
\draw[line width=0.12pt, black!45] (6.960,-0.550) -- (6.874,-0.522);
\draw[line width=0.12pt, black!45] (7.245,-0.541) -- (7.317,-0.501);
\draw[line width=0.14pt, black!51] (7.349,-0.570) -- (7.434,-0.558);
\draw[line width=0.13pt, black!49] (7.461,-0.572) -- (7.539,-0.566);
\draw[line width=0.12pt, black!47] (7.811,-0.538) -- (7.715,-0.546);
\draw[line width=0.13pt, black!50] (7.906,-0.581) -- (7.801,-0.590);
\draw[line width=0.14pt, black!52] (7.994,-0.535) -- (7.883,-0.544);
\draw[line width=0.14pt, black!53] (8.123,-0.600) -- (8.009,-0.610);
\draw[line width=0.13pt, black!52] (8.260,-0.590) -- (8.151,-0.599);
\draw[line width=0.12pt, black!48] (8.380,-0.600) -- (8.282,-0.608);
\draw[line width=0.11pt, black!42] (8.483,-0.533) -- (8.400,-0.540);
\draw[line width=0.14pt, black!50] (8.769,-0.600) -- (8.851,-0.593);
\draw[line width=0.15pt, black!57] (8.915,-0.542) -- (9.016,-0.533);
\draw[line width=0.16pt, black!61] (9.060,-0.588) -- (9.172,-0.579);
\draw[line width=0.16pt, black!62] (9.137,-0.541) -- (9.251,-0.532);
\draw[line width=0.16pt, black!62] (9.238,-0.587) -- (9.352,-0.577);
\draw[line width=0.11pt, black!41] (6.585,-0.428) -- (6.510,-0.403);
\draw[line width=0.12pt, black!45] (6.711,-0.444) -- (6.626,-0.411);
\draw[line width=0.12pt, black!48] (6.825,-0.438) -- (6.739,-0.390);
\draw[line width=0.12pt, black!47] (6.965,-0.449) -- (6.898,-0.380);
\draw[line width=0.12pt, black!46] (7.044,-0.436) -- (7.006,-0.351);
\draw[line width=0.12pt, black!44] (7.144,-0.444) -- (7.154,-0.355);
\draw[line width=0.12pt, black!45] (7.259,-0.433) -- (7.308,-0.356);
\draw[line width=0.11pt, black!43] (7.413,-0.456) -- (7.478,-0.405);
\draw[line width=0.11pt, black!40] (7.677,-0.475) -- (7.599,-0.473);
\draw[line width=0.12pt, black!47] (7.793,-0.428) -- (7.697,-0.433);
\draw[line width=0.13pt, black!52] (7.955,-0.416) -- (7.845,-0.425);
\draw[line width=0.14pt, black!53] (8.013,-0.480) -- (7.901,-0.489);
\draw[line width=0.14pt, black!53] (8.170,-0.446) -- (8.057,-0.455);
\draw[line width=0.13pt, black!52] (8.227,-0.410) -- (8.117,-0.419);
\draw[line width=0.13pt, black!49] (8.317,-0.431) -- (8.214,-0.440);
\draw[line width=0.11pt, black!40] (8.509,-0.410) -- (8.431,-0.417);
\draw[line width=0.14pt, black!54] (8.823,-0.462) -- (8.915,-0.454);
\draw[line width=0.16pt, black!59] (8.949,-0.417) -- (9.055,-0.408);
\draw[line width=0.16pt, black!61] (9.015,-0.440) -- (9.125,-0.431);
\draw[line width=0.16pt, black!62] (9.122,-0.441) -- (9.236,-0.431);
\draw[line width=0.16pt, black!61] (9.248,-0.434) -- (9.361,-0.425);
\draw[line width=0.11pt, black!40] (6.585,-0.374) -- (6.513,-0.343);
\draw[line width=0.11pt, black!42] (6.669,-0.330) -- (6.598,-0.287);
\draw[line width=0.12pt, black!46] (6.831,-0.331) -- (6.763,-0.265);
\draw[line width=0.12pt, black!47] (6.917,-0.301) -- (6.862,-0.221);
\draw[line width=0.12pt, black!47] (7.070,-0.376) -- (7.048,-0.282);
\draw[line width=0.12pt, black!47] (7.126,-0.311) -- (7.122,-0.214);
\draw[line width=0.12pt, black!45] (7.281,-0.298) -- (7.306,-0.210);
\draw[line width=0.11pt, black!41] (7.398,-0.311) -- (7.426,-0.234);
\draw[line width=0.12pt, black!45] (7.728,-0.296) -- (7.636,-0.298);
\draw[line width=0.13pt, black!49] (7.842,-0.348) -- (7.740,-0.354);
\draw[line width=0.13pt, black!51] (7.927,-0.365) -- (7.819,-0.374);
\draw[line width=0.14pt, black!53] (7.992,-0.305) -- (7.880,-0.314);
\draw[line width=0.14pt, black!53] (8.113,-0.297) -- (7.999,-0.306);
\draw[line width=0.13pt, black!51] (8.264,-0.334) -- (8.157,-0.343);
\draw[line width=0.12pt, black!47] (8.359,-0.331) -- (8.262,-0.339);
\draw[line width=0.11pt, black!42] (8.472,-0.303) -- (8.390,-0.310);
\draw[line width=0.14pt, black!54] (8.819,-0.372) -- (8.911,-0.365);
\draw[line width=0.16pt, black!58] (8.924,-0.331) -- (9.029,-0.322);
\draw[line width=0.16pt, black!61] (9.044,-0.306) -- (9.157,-0.296);
\draw[line width=0.16pt, black!62] (9.132,-0.331) -- (9.246,-0.322);
\draw[line width=0.16pt, black!62] (9.203,-0.370) -- (9.317,-0.361);
\draw[line width=0.11pt, black!42] (6.699,-0.271) -- (6.634,-0.220);
\draw[line width=0.12pt, black!45] (6.841,-0.254) -- (6.784,-0.184);
\draw[line width=0.12pt, black!46] (6.949,-0.255) -- (6.905,-0.170);
\draw[line width=0.12pt, black!46] (7.059,-0.222) -- (7.037,-0.130);
\draw[line width=0.12pt, black!47] (7.118,-0.258) -- (7.110,-0.162);
\draw[line width=0.12pt, black!45] (7.222,-0.225) -- (7.228,-0.135);
\draw[line width=0.11pt, black!41] (7.354,-0.205) -- (7.356,-0.125);
\draw[line width=0.12pt, black!46] (7.735,-0.194) -- (7.640,-0.200);
\draw[line width=0.13pt, black!49] (7.812,-0.233) -- (7.711,-0.242);
\draw[line width=0.13pt, black!51] (7.902,-0.199) -- (7.794,-0.212);
\draw[line width=0.14pt, black!53] (8.060,-0.231) -- (7.946,-0.244);
\draw[line width=0.14pt, black!53] (8.155,-0.196) -- (8.042,-0.210);
\draw[line width=0.13pt, black!51] (8.248,-0.263) -- (8.141,-0.273);
\draw[line width=0.12pt, black!47] (8.358,-0.254) -- (8.261,-0.263);
\draw[line width=0.11pt, black!42] (8.469,-0.256) -- (8.386,-0.263);
\draw[line width=0.14pt, black!51] (8.761,-0.239) -- (8.847,-0.231);
\draw[line width=0.15pt, black!57] (8.877,-0.191) -- (8.978,-0.182);
\draw[line width=0.16pt, black!60] (8.978,-0.230) -- (9.087,-0.221);
\draw[line width=0.16pt, black!62] (9.169,-0.212) -- (9.283,-0.203);
\draw[line width=0.16pt, black!61] (9.226,-0.232) -- (9.339,-0.223);
\draw[line width=0.11pt, black!40] (6.840,-0.102) -- (6.799,-0.035);
\draw[line width=0.11pt, black!42] (6.928,-0.111) -- (6.894,-0.036);
\draw[line width=0.11pt, black!42] (7.072,-0.104) -- (7.051,-0.023);
\draw[line width=0.11pt, black!43] (7.159,-0.138) -- (7.148,-0.052);
\draw[line width=0.11pt, black!42] (7.228,-0.130) -- (7.219,-0.048);
\draw[line width=0.11pt, black!40] (7.514,-0.097) -- (7.439,-0.076);
\draw[line width=0.11pt, black!43] (7.627,-0.115) -- (7.540,-0.114);
\draw[line width=0.12pt, black!47] (7.732,-0.080) -- (7.635,-0.089);
\draw[line width=0.13pt, black!51] (7.847,-0.085) -- (7.740,-0.097);
\draw[line width=0.13pt, black!52] (7.894,-0.103) -- (7.785,-0.115);
\draw[line width=0.14pt, black!53] (7.999,-0.091) -- (7.886,-0.105);
\draw[line width=0.14pt, black!53] (8.111,-0.082) -- (7.998,-0.096);
\draw[line width=0.13pt, black!51] (8.251,-0.122) -- (8.145,-0.135);
\draw[line width=0.12pt, black!45] (8.394,-0.102) -- (8.304,-0.114);
\draw[line width=0.11pt, black!40] (8.490,-0.140) -- (8.413,-0.150);
\draw[line width=0.13pt, black!50] (8.732,-0.132) -- (8.814,-0.122);
\draw[line width=0.14pt, black!53] (8.790,-0.100) -- (8.881,-0.089);
\draw[line width=0.15pt, black!58] (8.909,-0.086) -- (9.014,-0.073);
\draw[line width=0.16pt, black!61] (9.034,-0.141) -- (9.146,-0.129);
\draw[line width=0.16pt, black!62] (9.096,-0.157) -- (9.210,-0.146);
\draw[line width=0.16pt, black!61] (9.203,-0.108) -- (9.316,-0.094);
% the ark's road, Cariathiarim to Jerusalem (2 Kings 6)
\draw[line width=0.5pt, dash pattern=on 2.6pt off 1.8pt, -{Stealth[length=1.6mm]}]
  (5.85,-0.89) .. controls (6.6,-1.03) and (7.6,-1.07) .. (8.35,-0.95);
\fill (5.65,-0.83) circle (0.06); \node[stiny, anchor=south, fill=white, inner sep=0.8pt] at (5.72,-0.71) {Cariathiarim};
\draw[line width=0.5pt] (7.0,-1.33) circle (0.09);
\node[stiny, anchor=north] at (7.0,-1.49) {Ein Karem (a tradition, six centuries later)};
\fill (8.55,-0.97) circle (0.07); \draw[line width=0.4pt] (8.55,-0.97) circle (0.12);
\node[stiny, anchor=north] at (8.6,-1.13) {\scshape Jerusalem};
\end{tikzpicture}
\sketchcaption{The Visitation laid over the ark's procession, at the honest
weight of each contact; the wording on both sides is the Douay's, quoted
whole in the text. Below, the one patch of ground both stories cross. Luke
names no town, and the drawing keeps Ein Karem an open circle: a venerable
tradition, first attested by a pilgrim's distance six centuries on.}
\end{sketchfig}


\subsection{The overshadowing}

``The Holy Ghost shall come upon thee, and the power of the most High shall
overshadow thee'' (Lk~1:35). The verb beneath ``overshadow'' is
\term{episkiazein}, and its Old Greek career is short: exactly four
occurrences in the whole Septuagint. Twice it is the Psalter's protective
shadow, once wealth's glory in Proverbs. The fourth---the one
theologically laden use---is the cloud at the completion of the
tabernacle: Moses could not enter the tent of meeting, because the cloud
overshadowed it (\term{epeskiazen ep' autēn hē nephelē}) and the
tabernacle was filled with the glory of the Lord (Ex~40:35 LXX). The link
lives only in the Greek. Douay Exodus reads ``The cloud covered the
tabernacle\ldots\ the cloud covering all things'' (Ex~40:32--33), while
the English word ``overshadow'' surfaces at the New Testament end, at
Lk~1:35 and again at Heb~9:5, ``the cherubims of glory overshadowing the
propitiatory.''\endnote{Douay Ex~40:32--33 read at
\url{https://drbo.org/chapter/02040.htm}; the pericope stands at 40:34--35
in RSV and Hebrew-numbered editions, and at 40:28--29 in the Greek church
text's numbering (Rahlfs numbering is used in this study). LXX
\term{episkiazō} occurrences: Ex~40:35; Ps~90:4; Ps~139:8; Prov~18:11. New
Testament: Mt~17:5; Mk~9:7; Lk~1:35; Lk~9:34; Acts~5:15. Num~9:18--22,
sometimes cited in popular treatments as further occurrences, in fact uses
the simple \term{skiazō} (and once \term{skepazō}), not the compound.} And Luke
himself tells us what the verb means to him. He uses it again at the
transfiguration, where a cloud came and overshadowed the disciples
(Lk~9:34, \term{egeneto nephelē kai epeskiazen autous})---the theophanic
glory-cloud, not mere shelter. That internal control is the strongest
single piece of evidence that Lk~1:35 intends Shekinah imagery.

Now the precisions, both of them. First, the object of the overshadowing
in Exodus is the tent of meeting, not the ark as such: the ark stands
within the tent (Ex~40:20--21). Strictly taken, the figure Lk~1:35
supports is Mary as the overshadowed tabernacle, the place of the
glory-cloud; the ark is one step further in. Second, that step is shorter
than it once appeared. The adversarial re-run of the whole shadow-verb
family found the simple verb \term{skiazein} sitting directly on the ark
itself. At 1~Paralipomenon~28:18 LXX the cherubim spread their wings
\term{skiazontōn epi tēs kibōtou}, overshadowing the ark of the covenant
of the Lord---Douay: ``spreading their wings, and covering the ark''---and
at Ex~38:8 LXX the two cherubim stand \term{skiazonta} over the
propitiatory, with the veil that shadows the ark for the march and the
cloud that shall shadow Mount Sion in the same verb-family.\endnote{1~Par~28:18 LXX
verified in the corpus and at
\url{https://www.blueletterbible.org/lxx/1ch/28/18/}; Douay wording at
\url{https://drbo.org/chapter/13028.htm}. Ex~38:8 LXX (cherubim
\term{skiazonta} over the \term{hilastērion}; the Hebrew ordering of the
chapter differs); Num~4:5 (the veil \term{to syskiazon} with which the ark
is covered when camp breaks); Isa~4:5 (\term{skiasei nephelē} over Zion).}
The shadow-vocabulary of the Greek Bible, in short, is sanctuary furniture
vocabulary---cherub wings over the ark, cloud over the tent---though Luke's
own compound \term{episkiazein} never rests on the ark anywhere in the
Septuagint. Both facts belong in print.

The soberest critic of this parallel, Marie Isaacs, granted the verb and
denied the inference: ``it cannot be assumed that Luke is employing a
typology whereby Mary represents the tabernacle in whom the Shekinah
resides''. Her stated ground was that the Shekinah was not confined to
the tabernacle or the ark, with the Psalter's protective usage as her
supporting observation.\endnote{Marie
E.~Isaacs, ``Mary in the Lucan Infancy Narrative,'' \work{The Way
Supplement} 25 (1975), \url{https://www.theway.org.uk/back/s025Isaacs.pdf},
read in full for this study. The quoted sentence and her stated ground
(``the Shekinah was not confined to the tabernacle or the ark'') are
verbatim from the article; she observes that the verb occurs elsewhere
``similarly to signify the divine presence'' (Ps~90:4; 139:8 LXX).} That is
fairly said, and it does not touch Lk~9:34. A rare verb, an uncontested
glory-cloud sense in Luke's own usage, a tent rather than an ark at the
Exodus end: that is the exact weight of the first echo.

\subsection{The question at the door}

When the ark set out for Jerusalem and Oza lay dead, David was afraid,
and asked: ``How shall the ark of the Lord come to me?''
(2~Kings~6:9; LXX \term{pōs eiseleusetai pros me hē kibōtos kyriou?}).
Elizabeth, filled with the Holy Ghost, asks: ``And whence is this to me,
that the mother of my Lord should come to me?'' (Lk~1:43; \term{pothen moi
touto hina elthē hē mētēr tou kyriou mou pros eme?}). Shared: a rhetorical
question asked of oneself; a sacred bearer coming---\term{pros me}, ``to
me''---and ``the Lord'' in the genitive at the heart of both phrases. Not
shared: the interrogative itself (\term{pōs} against \term{pothen moi
touto}) and the register, David's dread against Elizabeth's joy.

The critics have a rival text, and it deserves the body of the page, not
a footnote. Raymond Brown observed that when David comes to
Araunah's threshing floor---the future temple site---Araunah asks, in
Brown's rendering, ``What is this, that my lord the king has come to his
servant?'' (2~Kings~24:21; LXX \term{ti hoti ēlthen ho kyrios mou ho
basileus pros ton doulon autou?}). ``This question also resembles Elizabeth's question,'' Brown
concludes, ``and it does not concern the Ark''; on his reading ``the
connecting link in the Lucan reminiscences may be David rather than the
Ark.''\endnote{Raymond E.~Brown, \work{The Birth of the Messiah} (updated
ed.; New York: Doubleday, 1993), pp.~344--345, quoted here as transmitted
verbatim in Kozłowski 2018, n.~9 (see below); Brown's pages were not
opened at first hand for this study. His dedicated subsection ``The Ark of
the Covenant in 1:35?''\ stands at p.~327 of the same work, as reported
from its table of contents.} The rejoinder, pressed by Jan
Kozłowski, is about content rather than syntax. Araunah's question is
court courtesy to a superior. David's question and Elizabeth's are the
only two shaped as numinous unworthiness---the creature asking how the
bearer of the Lord's presence can enter a private life.\endnote{Jan M.~Kozłowski, ``Mary
as the Ark of the Covenant in the Scene of the Visitation (Luke 1:39--56)
Reconsidered,'' \work{Warszawskie Studia Teologiczne} 31/1 (2018)
108--116, \url{https://czasopismowst.pl/index.php/wst/article/view/159},
read in full for this study; his supporting argument on Lk~1:42 is
``\,`The Fruit of Your Womb' (Luke 1,42) as `The Lord God, Creator of
Heaven and Earth' (Judith 13,18),'' \work{Ephemerides Theologicae
Lovanienses} 93 (2017) 339--342, cited here as reported. A popular
restatement is his ``Biblical Intertextuality: The Virgin Mary as the Ark
of the Covenant,'' \work{Antigone Journal}, June 2022.} The honest verdict:
a genuine motif-parallel with partial verbal support, an allusion at
most, never a quotation---and at this single point the debate is
genuinely open.

\subsection{Elizabeth's cry}

The philologically strongest item in the cluster is a single word. Luke
writes that Elizabeth \term{anephōnēsen kraugē megalē}---cried out with a
loud cry---and blessed Mary among women and the fruit of her womb
(Lk~1:42). The verb \term{anaphōnein} occurs nowhere else in the New
Testament. In the whole Septuagint it occurs exactly five times, every
one in Paralipomenon, every one a Levite's liturgical acclamation, four
of the five in ark ceremonies: the procession that brings the ark up to
Jerusalem, the ministers appointed before the ark, the cultic music of
Heman and Idithun (the one occurrence with no ark in the verse), and the
ark's installation in Solomon's temple. That last acclamation ends with
the house filled with the cloud of the glory of the Lord---the same
Exodus~40 cluster that stands behind Lk~1:35.\endnote{LXX \term{anaphōneō}:
1~Par~15:28; 16:4, 5, 42; 2~Par~5:13---the five occurrences confirmed by
independent whole-corpus search; Kozłowski (2018, pp.~113--114) states the
proportion as eighty percent. 2~Par~5:13 ends \term{kai ho oikos eneplēsthē
nephelēs doxēs kyriou}. The critical text of Lk~1:42 reads \term{kraugē
megalē} (so SBLGNT, Nestle 1904, Westcott-Hort); the Byzantine and Received
texts read \term{phōnē megalē}---the verb itself is stable across all
traditions.} There is a further convergence: the ark goes up to Jerusalem
\term{meta kraugēs}, with shouting (2~Kingdoms~6:15), and the critical
text of Lk~1:42 gives Elizabeth precisely a \term{kraugē}. And Elizabeth is
a daughter of Aaron (Lk~1:5), which gives the Levitical color of the verb
somewhere to rest.

The caveats are printed with the claim. Five occurrences are a small base
on which to declare a technical term. In Greek at large \term{anaphōnein}
is an ordinary verb of crying out, so the Septuagintal distribution is a
fact about a corpus, not a property of the word. And a hapax proves that
Luke chose a rare word, not why he chose it. The statistic is suggestive,
verified, and not a demonstration. That is its weight.

\subsection{Three months and a blessed house}

The record of the first procession says that ``the ark of the Lord abode
in the house of Obededom the Gethite three months: and the Lord blessed
Obededom, and all his household''
(2~Kings~6:11). ``And Mary abode with her about three months; and she
returned to her own house'' (Lk~1:56). The Greek match is exact where it
matters: \term{mēnas treis}, ``three months,'' verbatim on both sides,
Luke adding only ``about''. The Douay English adds a small bonus of its
own, the same verb ``abode'' in both verses. The blessing motif stands
beside the number: the Lord blessed (\term{eulogēsen}) all Obededom's
house because of the ark, and Elizabeth pronounces Mary blessed and the
fruit of her womb blessed (\term{eulogēmenē\ldots\ eulogēmenos},
Lk~1:42)---though Luke never says Zachary's house was blessed, and
honesty keeps that difference visible.

Two deflationary facts are also visible, and the study prints them. The
three months are narratively motivated: Mary arrives in Elizabeth's sixth
month (Lk~1:36) and stays until John's birth is due; no typology is needed
to generate the number. And the Book of Judith stands nearby. Israel
rejoiced for three months after Judith's victory (Jdt~16:20), and
Elizabeth's actual sentence of blessing demonstrably tracks Judith's
blessing---blessed art thou, daughter, by the most high God above all
women (Jdt~13:18: \term{eulogētē sy, thygater, tō theō tō hypsistō para
pasas tas gynaikas})---more closely than any ark text does.\endnote{Jdt
16:20 LXX: the people rejoicing in Jerusalem \term{epi mēnas treis}, Judith
remaining with them; Jdt~13:18 as transliterated in the text. Both loci
verified in the corpus for this study; the English glosses of the Judith
verses are editorial renderings of the Greek, not quotations of any
version.}

But the Judith echoes are polyphony, not rivals. Luke's manner in these
chapters is to layer scriptural voices---Hannah under the Magnificat,
Daniel under Gabriel---and a Judith harmony does not silence an ark
harmony. Indeed Kozłowski's revival of the whole argument begins from the
Judith verse, as we shall see. His judgment on this item is that
2~Kings~6:11 remains the closer parallel to Lk~1:56: in both, a
God-bearing presence abides three months in a private Judean house, and
blessing attends the stay.

\subsection{The leap: an honest negative}

Here the study prints a negative finding, because the popular form of this
parallel claims more than the texts give. John ``leaped'' in Elizabeth's
womb---\term{eskirtēsen} (Lk~1:41, 44)---and David danced before the ark;
the apologetic commonplace says they share a verb. They do not. In
2~Kingdoms~6 the Septuagint's David plays (\term{paizontes}, v.~5), strikes
up on instruments (\term{anekroueto}, v.~14---the Greek has actually
softened the Hebrew's whirling), and is seen dancing and playing
(\term{orchoumenon kai anakrouomenon}, v.~16). The verb \term{skirtan}
never occurs in the chapter. Its true Septuagintal partner is Gen~25:22:
\term{eskirtōn de ta paidia en autē}---the twins leaping in Rebekah's
womb, in a verse whose question, \term{hina ti moi touto?}, is formally as
close to Lk~1:43 as anything in Kingdoms. Luke may well be interlacing a
second typology here---children leaping in the womb before the birth of
Israel, and of the new Israel---and Isaacs made the point with relish:
``Even the words for `dance' and `leap' are not the same in the
greek.''\endnote{LXX \term{skirtaō}: Gen~25:22; Ps~113:4, 6; Wis~17:18
(Rahlfs numbering); Joel~1:17; Mal~3:20; Jer~27:11 LXX (= MT 50:11)---seven
occurrences, none in 2~Kingdoms~6; in the New Testament the verb is
exclusively Lucan (Lk~1:41, 44; 6:23). Isaacs's sentence is quoted as
printed (lowercase ``greek''), the Gen~25:22 observation credited by her to
Goulder and Sanderson, \work{JTS} 8 (1957), p.~21. Laurentin himself
glossed \term{skirtan} from Mal~3:20 and Ps~113:4--6 and a Wisdom locus
transmitted as ``17,9''---a phantom as printed, since the only Wisdom
occurrence is 17:18 (17:19 in some numberings); and his fifth point lists
the ``bonds joyeux de David et de Jean-Baptiste'' among the shared
manifestations of joy (as quoted in Kozłowski 2018, n.~8). The
motif-parallel is thus Laurentin's own; the same-verb claim is not his.}

Two precisions keep the negative honest in both directions. Laurentin
never claimed the same verb. He glossed John's leap from Malachias and
the Psalter---leaping at the coming of the Lord---and listed the joyful
bounds of David and of John together as a shared motif; the lexical
equation is a popular overreach beyond what he wrote. And the motif
equation itself has an ancient, if late, witness. In the second century after Christ the
Jewish translator Symmachus rendered David's whirling at 2~Kings~6:16 with
precisely this verb: \term{skirtōnta kai kanchazonta}, leaping and
laughing---the one ancient Greek Bible text in which David does
\term{skirtan} before the ark.\endnote{Field, \work{Origenis Hexaplorum
quae supersunt}, vol.~I (Oxford: Clarendon, 1875), p.~555, with n.~28 (the
reading rests on Codex Regius and Cod.~243, with the orthographic variant
\term{kakchazonta} and a misprint in Montfaucon); page image read at
\url{https://archive.org/details/origenhexapla01unknuoft}. Field prints
the present participle \term{skirtōnta}; the aorist sometimes quoted in
apologetic use exists nowhere. Symmachus worked in the later second
century A.D., after Luke.} Symmachus is post-Lucan; he cannot be what Luke
alluded to, and the verdict for Luke's own composition stands: no verbal
link. What Symmachus shows is smaller and still real: an ancient
translator, independent of the Church, found ``leaping'' the natural
Greek for what David did before the ark. The equation of the two
scenes is not a modern contrivance. It is a hexaplaric footnote, not a
proof.

\subsection{The hill country}

``Arose and went'': the formula \term{anastasa\ldots\ eporeuthē} (Lk~1:39),
matching \term{anestē kai eporeuthē Dauid} (2~Kingdoms~6:2), is
Septuagintal boilerplate---Abraham rises and goes to Moriah in the same
words (Gen~22:3), and dozens of journeys begin so. By itself it carries no
allusive weight, and the study assigns it none. What carries weight is the
cluster riding on it: arise, go, Juda, the ascent into hill country, a
destination house, a sacred burden.

Here the Greek Bible offers one curious texture. The Septuagint of
2~Kingdoms~6:2 has lost the place-name. Where the Hebrew has David set
out from Baale-Juda (that is, Kiriath-jearim), the Greek translator read
the noun as a common one,
producing \term{apo tōn archontōn Iouda}---from the rulers of Juda, on the
ascent (\term{en anabasei})---and the Vulgate, followed by the Douay's
``of the men of Juda,'' lost the place-name the same way. A reader of the
Greek or Latin Bible thus meets in 2~Kings~6:2 an unnamed Judahite starting
point and an ascent---a texture oddly congenial to Luke's unnamed ``city
of Juda'' in the hill country, though Chronicles keeps the name
(1~Par~13:5--6). The ark comes into the house of Obededom
(\term{eis oikon Abeddara}, 6:10--11); Mary enters the house of
Zachary (\term{eisēlthen eis ton oikon Zachariou}, Lk~1:40).

Luke names no town, and the tradition that names one must be dated as
tradition. The identification of the Visitation with Ein Karem, west of
Jerusalem, rests at its oldest on the pilgrim itinerary of Theodosius,
composed between about 518 and 530. It gives a distance only:
\term{De Hierusalem usque ubi habitavit sancta Helisabeth mater domini
Ioannis Baptistae milia V}---from Jerusalem to where St.~Elizabeth, the
mother of John the Baptist, lived, five miles---and no name; one
manuscript adds that holy Mary greeted Elizabeth there. The name itself
appears in the later liturgical tradition of Byzantine Jerusalem, as
reported.\endnote{Theodosius,
\work{De situ terrae sanctae} \S 35, in J.~Gildemeister's 1882 edition;
the manuscript S addition \term{salutavit sancta Maria Elisabeth} stands
in the apparatus; the English renderings in the text are editorial, not a
translation. Read at
\url{https://archive.org/details/theodosiusdesit00theogoog}; the
paragraph-numbering is Gildemeister's, and the CCSL~175 numbering has not
been confirmed for this study. The later attestation of the name
(Enqarim/Ain Karim) in the Byzantine Jerusalem liturgical tradition is
reported at second hand. Ein Karem lies some eight kilometers
west-southwest of Jerusalem's Old City and Kiriath-jearim some thirteen
kilometers west, as the standard atlases report.} Both Ein Karem and
Kiriath-jearim lie in the Judean hill country, and the popular claim that
Mary ``took the same road as the ark'' has no textual basis whatever: Luke
names neither town nor route. The tradition is venerable. It is not
evidence, and this study does not spend it as evidence.

\subsection{Daughter Zion, the companion typology}

One more Old Testament pattern stands over Luke~1, and it must be kept
distinct before it can be related. ``Hail, full of grace, the Lord is with
thee: blessed art thou among women'' (Lk~1:28, the Douay printing the
longer traditional text of the verse). The greeting \term{chaire},
followed by the announcement that the Lord is coming to dwell, is read by
a whole line of scholarship as the Septuagint's eschatological summons to
the Daughter of Zion: \term{chaire sphodra, thygater Siōn}---rejoice
greatly, daughter Sion (Soph~3:14)---with the King of Israel, the Lord, in
the midst of thee (3:15) and the Lord thy God in thee (3:17); so too
Zach~9:9. On this reading Gabriel greets Mary as the personified Zion
whom the Lord enters, and \term{kecharitōmenē}---a rare verb---names the
grace of that election.

The minimalist answer is verbal. The only word Lk~1:28 exactly shares
with Sophonias is \term{chaire} itself, which is also the everyday Greek
greeting; ``the Lord is with thee'' belongs to the call-formula spoken to
Gedeon and by Booz (Judg~6:12; Ruth~2:4), not to the Zion oracles. Isaacs
judged the verbal parallels ``are, in fact, slight,'' the womb-for-midst
equation ``wholly far-fetched''---while allowing that the
Daughter-of-Zion personification ``may have influenced Luke.''\endnote{Soph~3:14--17; Zach~9:9; Joel~2:21; Judg~6:12 (Codex
Vaticanus text; Alexandrinus varies in the epithet); Ruth~2:4---all
verified in the corpus. Douay Lk~1:28 prints the ``blessed art thou among
women'' clause with the Vulgate and Byzantine tradition; the critical
text carries it only at Lk~1:42. Isaacs's phrases as in her article
(above); Fitzmyer's treatment of \term{chaire} as an ordinary greeting is
reported at second hand. The Daughter-Zion exegesis descends from
S.~Lyonnet, ``Chaire kecharitōmenē,'' \work{Biblica} 20 (1939) 131--141
(reported), through Laurentin to Joseph Ratzinger, \work{Daughter Zion}
(San Francisco: Ignatius, 1983; German 1977), whose sentences---the
salutation ``modeled closely on Zephaniah 3:14-17''; Mary ``in person the
true Zion''; the praise of him whom she bears, ``the true, incorruptible,
indestructible Ark of the Covenant''---are quoted from the excerpt at
\url{https://www.catholicworldreport.com/2025/03/25/joseph-ratzinger-the-mystery-of-the-annunciation-is-the-mystery-of-grace/};
page numbers not established. Note the shape of Ratzinger's claim: the
title ``Ark'' passes finally to Christ, Mary being ark as his bearer.}

The point of printing this here is precision. Daughter Zion and the ark
are two different figures: the first makes Mary the city and people whom
the Lord indwells, the second the vessel that carries the Presence from
place to place. The Church's own prayer has fused them. The Catechism, in
its commentary on the Ave Maria, calls Mary the daughter of Zion in
person, the ark of the covenant, the place where the glory of the Lord
dwells (CCC 2676); Ratzinger fuses them deliberately, handing the ark
title finally to Christ with Mary as bearer. That fusion is an act of
reception, and its story is told in section~32. Here it is enough to say
what the texts show: two distinct patterns, each with its own verbal
skeleton, laid over the same seventeen sentences.

\subsection{The debate as a debate}

The modern argument has a genealogy. Eric Burrows first developed the
tabernacle-and-ark reading of Luke~1 in essays published in 1940;
Stanislas Lyonnet argued the Daughter-Zion reading of the greeting in
1939. René Laurentin built both into the classic statement,
\work{Structure et théologie de Luc I--II} (1957),
where the Visitation's chosen traits refer the reader ``de façon
convergente''---convergently---to the narrative of the ark's transfer, and
Mary is ``le type de l'arche d'alliance.''\endnote{E.~Burrows, \work{The
Gospel of the Infancy and Other Biblical Essays}, ed.~E.~F. Sutcliffe
(London: Burns Oates, 1940), esp.~pp.~47f., reported via Isaacs's notes
(Laurentin p.~191: ``following Burrows, op.~cit., pp 47f''); Lyonnet as in
the preceding note; R.~Laurentin, \work{Structure et théologie de Luc
I--II} (Paris: Gabalda, 1957; repr.~1964)---Daughter Zion pp.~64--71, the
overshadowing p.~73, the ark parallels pp.~79--81, ``le type de l'arche
d'alliance'' p.~80; his summary sentence at p.~79, ``Le récit --- très
stylisé --- est composé de traits choisis qui nous renvoient de façon
convergente au récit du transfert de l'arche à Jérusalem par David,'' is
quoted as transmitted in Kozłowski 2018 (whose numbering of Laurentin's
five points is Kozłowski's own); the 1957/1964 pagination equivalence has
not been confirmed at first hand.} The argument's form is the one
weighed at the toll-gate: cumulative and convergent. No single echo is
claimed as proof. The claim is that so many independent contacts within
the eighteen verses of one pericope (Lk~1:39--56) are unlikely to be
accident.

The critics answer by dismantling the parallels one at a time and offering
a rival frame. Brown's Araunah text has been quoted above; his positive
proposal is that the reminiscences are Davidic rather than
arkological---Mary enters the story David entered, not the ark's. Fitzmyer:
``But this is subtle. If, indeed, the story may be compared with 2~Sam
24:21, then what connection does it have with the ark?'' Schürmann held
that the ark narrative aims at Jerusalem's election, so that Mary shares
with the ark only the direction of travel, never the destination, and that
the fear-laden three months of the ark are ``sehr unpassend''---very
unfittingly---compared with Mary's stay. Bovon was briefest: ``Die
Symbolik Maria = Lade ist zu weit hergeholt''---the symbolism is too
far-fetched. Strauss found it ``unlikely\ldots\ that Luke would be drawing
so subtle an allusion here''; Isaacs, that the case ``rests upon the
flimsiest of correspondences,'' and that ``it is to go beyond the evidence
to suggest that she is a type of the ark of God.''\endnote{J.~A.~Fitzmyer,
\work{The Gospel According to Luke I--IX} (Anchor Bible 28; New York:
Doubleday, 1981), p.~364; H.~Schürmann, \work{Das Lukasevangelium~I}
(HThKNT; Freiburg: Herder, 1969), pp.~64--65 n.~161 (``Das furchterfüllte
dreimonatige Belassen der Lade\ldots\ wird sehr unpassend\ldots\
verglichen''; Mary has ``nur `la direction', keineswegs aber das
Reiseziel''); F.~Bovon, \work{Das Evangelium nach Lukas~I} (EKK, 1989),
p.~86 n.~43, = \work{Luke~1} (Hermeneia, 2002), p.~59 n.~46, with his
\work{Luke the Theologian}, pp.~183--184, against the identification of
the child with YHWH that underlies the recent revival; M.~L.~Strauss,
\work{The Davidic Messiah in Luke-Acts} (JSNTSup 110; Sheffield, 1995),
p.~96; H.~Rice, \work{Behold, Your House Is Left to You} (Pickwick,
2016), p.~63 (``Less persuasively, Laurentin connects\ldots''). All of
these are quoted as transmitted verbatim in Kozłowski 2018, nn.~9--10,
whose quotations were checked in full for this study; the originals were
not opened at first hand. Isaacs is quoted from her article directly.}
This is a school with premises, and the study reports the premises with
the verdicts: on strict verbal-echo criteria, an allusion must be lexically
demonstrable---the same words, not the same shape---or it is not there; and
liturgical reception is not admitted as evidence of what Luke meant.

The argument did not end in the commentaries. Kozłowski's 2017--2018
intervention re-armed Laurentin from an unexpected side. Lk~1:42 blesses
Mary and the fruit of her womb in the syntax Judith~13:18 uses to bless a
woman and, in the phrase of Kozłowski's own title, ``The Lord God, Creator
of Heaven and Earth''---so the child in the womb is acclaimed in Creator's
language, and if the burden is the Lord, the bearer is what the ark was.
To this he added the corrected concordance facts for \term{anaphōnein}
set out above. From the Reformed side, Jacobus de Koning argued the
typology in 2020 from 2~Kings~6 with the Chronicles acclamation texts. He
noted that it is ``grootliks afwesig''---largely absent---from the
mainstream commentaries, and proposed it as a challenge for Protestant
rethinking of Mariology.\endnote{Jacobus d.~W. de Koning, ``Maria as anti-tipe van die
verbondsark in Lukas 1:39--56: 'n Uitdaging tot Protestantse herbesinning
van die Marialogie,'' \work{In die Skriflig/In Luce Verbi} 54(1), a2561
(17 June 2020),
\url{https://indieskriflig.org.za/index.php/skriflig/article/view/2561},
read in full for this study; the ``grootliks afwesig'' judgment on the
commentary tradition (Bock, Fitzmyer, Green) stands at his p.~10.} And behind
both stands an identifiable Catholic school with institutional addresses:
Aristide Serra at the Marianum, whose Old Testament prefigurations of
Mary include ``the tent of the ark of the covenant''; André Feuillet at
the Institut Catholique de Paris; Ignace de la Potterie at the Biblicum,
for whom the Daughter-of-Zion promises are realized in a definite woman.
This study cites their loci as reported, not as texts opened at first
hand.\endnote{A.~Serra, \work{La Donna dell'Alleanza:
Prefigurazioni di Maria nell'Antico Testamento} (Padova: Messaggero,
2006), whose chapter plan includes ``La tenda dell'Arca dell'alleanza''
(publisher's table); A.~Feuillet, \work{Jésus et sa mère} (Paris: Gabalda,
1974; ET \work{Jesus and His Mother}, 1984), with his ``Le Messie et sa
Mère d'après le chapitre XII de l'Apocalypse,'' \work{Revue Biblique} 66
(1959) 55--86; I.~de la Potterie, \work{Mary in the Mystery of the
Covenant}, trans.~B.~Buby (New York: Alba House, 1992). Laurentin's
compact late statement of the ark argument stands at his \work{A Short
Treatise on the Virgin Mary} (ET 1991), pp.~27--30. All of these loci are
reported from the survey literature and publishers' records assembled in
the study's research files; none was read at first hand this session.}

Name the divide honestly and it names itself: methodological, not
confessional. The thesis was born among Jesuits; its sharpest critics are
Catholic priests in good standing; its newest defender is Reformed. What
separates the parties is what counts as an allusion---strict lexical
identity, or convergent pattern---and whether the reading Church's use of
a text is evidence about the text or only about the Church.

\subsection{What the philology honestly yields}

Count the inventory. One rare verb whose only laden Septuagintal use is
the glory-cloud over the completed tabernacle, re-used by Luke himself for
the transfiguration cloud---with the tent, not the ark, beneath it. One
exact two-word match, ``three months,'' in the same structural position,
with a narrative explanation available and a Judith harmony beside it. One
New Testament hapax whose entire Old Greek career is Levitical acclamation,
four-fifths of it around the ark---five occurrences in all. One question
shaped like David's question, with a rival Davidic text that lacks the
ark. A journey texture congenial to the parallel but proving nothing
alone. One honest negative, where the popular claim fails and only a
post-Lucan translator supplies the missing verb. And one companion
typology, Daughter Zion, distinct in figure, fused with the ark only
downstream, in the Church's prayer.

The verdict this study prints is the one the evidence will bear: the
Visitation cluster is a convergent set of which no single strand compels.
A reader who requires lexical demonstration, echo by echo, will judge with
Brown and Bovon that the allusion is unproven---and on that criterion the
judgment is correct. A reader who admits cumulative design in a writer as
deliberate as Luke, and who weighs the improbability of five independent
contacts in eighteen verses, will judge with Laurentin's school that the
evangelist knew what procession he was describing---and no verified fact
above forbids it. Two verdicts, two premises; the rule paid at the
toll-gate holds the space between them: attestation is not reception, and
the absence of either is not falsity. The philology is what it is. What
the Church has made of the reading is a separate question with separate
evidence, and it occupies the next two sections.

One thing more, and it changes the frame. The whole debate reported above
was conducted, on both sides, as if the reading were born with Burrows in
1940. It was not. Around the year 500, a Syriac poet was already preaching
the Visitation with the ark---Mary an ark full of mysteries, John's leap
answering David's dance before it. His name was Jacob of Serugh, and the
Fathers are the next section's business. The philology can carry the
reader no further; the Fathers take him up.
